\chapter{System Design}

\section{Introduction}

This chapter presents the comprehensive system design for the E-Commerce Microservices Platform, translating the requirements specified in the previous chapter into concrete architectural decisions, component designs, data models, API specifications, and integration patterns. The design process follows established software engineering principles including Domain-Driven Design for service boundary definition, microservices architecture patterns for distributed system construction, event-driven architecture for asynchronous communication, and cloud-native design principles for scalability and resilience.

The system design addresses multiple abstraction levels ranging from high-level architectural decomposition showing service boundaries and communication patterns, through detailed component designs showing internal service structure and class relationships, to data models defining database schemas and event schemas specifying message formats. The design decisions documented in this chapter are justified based on requirements satisfaction, industry best practices, technology capabilities, and trade-off analysis considering factors such as complexity, performance, maintainability, and operational overhead.

\section{Architectural Design}

The architectural design establishes the fundamental structure of the system, defining the major components, their responsibilities, interactions, and the principles governing their organization. The architecture follows microservices patterns with event-driven communication, implementing a layered approach where infrastructure services provide foundational capabilities, business services implement core e-commerce functionality, and the frontend application provides the user interface.

\subsection{High-Level Architecture}

The platform architecture comprises three distinct layers: the presentation layer consisting of the React-based single-page application that provides the user interface for customers and administrators, the application layer consisting of eight microservices implementing business logic and data management, and the infrastructure layer consisting of databases, message brokers, caches, service discovery, and monitoring systems that support the application layer.

The React frontend application communicates exclusively with the API Gateway, which serves as the single entry point for all client requests. The API Gateway performs request routing based on URL patterns, forwarding requests to appropriate backend services: requests to the slash api slash users path are routed to the User Service on port eight thousand eight hundred and one, requests to slash api slash products are routed to the Product Service on port eight thousand eight hundred and two, requests to slash api slash orders are routed to the Order Service on port eight thousand eight hundred and three, requests to slash api slash payments are routed to the Payment Service on port eight thousand eight hundred and four, and requests to slash api slash inventory are routed to the Inventory Service on port eight thousand eight hundred and five. The Notification Service does not expose REST endpoints and is not directly accessible through the API Gateway, instead consuming events from Kafka to send notifications asynchronously.

Inter-service communication follows two patterns: synchronous REST API calls for request-response interactions where immediate responses are required, and asynchronous event-driven communication through Apache Kafka for loose coupling and eventual consistency. The Order Service orchestrates distributed transactions through the Saga pattern, publishing OrderCreated events to Kafka that trigger inventory reservation in the Inventory Service and payment processing in the Payment Service, with each service publishing response events that advance the Saga workflow.

The architecture implements the database-per-service pattern where each microservice owns its dedicated PostgreSQL database instance, preventing direct database access between services and enforcing API-based communication. This design provides strong service isolation, independent schema evolution, and technology flexibility, though it introduces complexity for queries spanning multiple services and requires distributed transaction patterns for data consistency.

\subsection{Service Decomposition Strategy}

Service decomposition follows Domain-Driven Design principles, identifying bounded contexts based on business capabilities, data ownership, and team organization. Each bounded context becomes a microservice with clear responsibilities, autonomous deployment, and independent scaling.

The User Service owns the user management bounded context, handling user registration, authentication, profile management, and role-based access control. This service owns user identity data, credential storage, and authorization information, providing authentication services to other services through JWT token validation and user information extraction. The service boundary aligns with the natural cohesion of identity management functionality and the need for centralized authentication logic.

The Product Service owns the product catalog bounded context, managing product information, categories, search functionality, and catalog browsing capabilities. This service owns all product-related data including names, descriptions, prices, categories, and inventory levels, providing query and search operations optimized for customer-facing product discovery. The service boundary reflects the distinct data model and access patterns of catalog management compared to other domains.

The Order Service owns the order management bounded context, orchestrating order creation, Saga workflow execution, order status tracking, and order history management. This service owns order data including order headers, line items, status history, and Saga state, coordinating with Inventory and Payment services through event-driven communication without directly accessing their databases. The service boundary encapsulates the complex order processing logic and Saga orchestration responsibilities.

The Payment Service owns the payment processing bounded context, handling payment gateway integration, transaction recording, refund processing, and payment status management. This service owns payment transaction data including payment records, gateway transaction identifiers, amounts, and statuses, integrating with external payment providers such as Stripe and PayPal through their respective APIs. The service boundary isolates payment processing complexity and PCI DSS compliance requirements from other services.

The Inventory Service owns the inventory management bounded context, managing stock levels, inventory reservations, stock adjustments, and low-stock alerts. This service owns inventory data including product stock quantities, reservation records, and stock movement history, consuming order events to manage reservations and publishing inventory events to acknowledge successful reservations or report failures. The service boundary separates inventory tracking logic from order processing logic while maintaining tight integration through events.

The Notification Service owns the notification delivery bounded context, consuming events from Kafka and sending notifications through email, SMS, and push notification channels. This service owns notification templates, delivery status records, and user notification preferences, providing reliable notification delivery with retry mechanisms for transient failures. The service boundary isolates notification delivery complexity from business logic services.

The API Gateway and Eureka Server serve as infrastructure services providing cross-cutting concerns rather than business functionality. The API Gateway implements the API Gateway pattern for request routing, authentication, rate limiting, and circuit breaking. The Eureka Server implements service discovery and registration, maintaining the service registry that enables dynamic service lookup and client-side load balancing.

\section{Component Design}

Component design specifies the internal structure of each microservice, including layer organization, class responsibilities, data access patterns, and integration mechanisms.

\subsection{Layered Architecture Pattern}

All microservices implement a consistent layered architecture pattern with four distinct layers: the presentation layer containing REST controllers that handle HTTP requests and responses, the service layer containing business logic implementations and orchestration, the repository layer containing data access abstractions using Spring Data JPA, and the domain layer containing entity classes representing database tables and value objects encapsulating business concepts.

The presentation layer implements RESTful API endpoints using Spring Web annotations including RequestMapping for endpoint path specification, GetMapping for retrieval operations, PostMapping for creation operations, PutMapping for update operations, and DeleteMapping for deletion operations. Controllers delegate business logic to service layer components, perform input validation using bean validation annotations, and return appropriate HTTP status codes including two hundred for successful operations, two hundred and one for resource creation, four hundred for validation errors, four hundred and one for authentication failures, four hundred and three for authorization failures, four hundred and four for resource not found, and five hundred for server errors. Controllers also handle exception translation, catching domain-specific exceptions and converting them to standardized error responses with appropriate status codes and error messages.

The service layer implements business logic, transaction management, and inter-service communication. Service classes are annotated with Spring's Service stereotype and Transactional annotations, with methods implementing business operations that may span multiple repository calls and external service invocations. The service layer enforces business rules, performs complex validations, orchestrates multi-step workflows, and handles error scenarios with appropriate compensating actions. For the Order Service, the service layer implements Saga orchestration logic, publishing events to Kafka and consuming response events to advance the order state machine.

The repository layer provides data access abstractions using Spring Data JPA repository interfaces that extend JpaRepository, inheriting standard CRUD operations and enabling custom query methods through method naming conventions or JPQL annotations. Repository interfaces specify entity types and primary key types, with Spring Data JPA automatically generating implementation classes at runtime. Custom queries leverage JPQL for complex retrievals including joins, aggregations, and subqueries, with pagination support through Pageable parameters and sorting through Sort parameters. The repository layer isolates data access logic from business logic, enabling database technology changes without affecting service layer code.

The domain layer contains entity classes annotated with JPA annotations including Entity to mark classes as persistent entities, Table to specify table names, Id to designate primary key fields, GeneratedValue for auto-generation strategies, Column for column specifications including names, nullability, and uniqueness constraints, and relationships including OneToMany, ManyToOne, and ManyToMany with cascade and fetch type configurations. Entity classes implement business methods that operate on entity state, encapsulating domain logic within the entities themselves following rich domain model principles rather than anemic domain models where entities contain only data without behavior.

\subsection{Order Service Saga Implementation}

The Order Service implements the most complex component design due to Saga orchestration responsibilities, requiring careful state machine design, event handling, and compensating transaction management.

The Saga state machine maintains order status transitions through well-defined states and transitions. The initial state is PENDING, entered when an order is first created. From PENDING state, the order can transition to CONFIRMED upon receiving StockReserved and PaymentCompleted events in sequence, or to CANCELLED upon receiving StockReservationFailed or PaymentFailed events, or upon receiving an explicit cancellation request from the customer. The state machine enforces valid transitions, rejecting invalid transitions such as attempting to cancel an order that has already been shipped.

Event handlers consume Kafka events using Spring Kafka's KafkaListener annotation, specifying topic names and consumer group identifiers. The Order Service subscribes to the inventory-events topic to receive StockReserved and StockReservationFailed events, and to the payment-events topic to receive PaymentCompleted and PaymentFailed events. Event handlers deserialize event payloads from JSON format into event objects, validate event structure and content, retrieve the associated order from the database, validate the current order status allows the requested transition, update the order status, persist the status change, and publish subsequent events if required.

Compensating transaction handlers execute when Saga steps fail, reversing the effects of previously completed steps to maintain consistency. When payment processing fails after successful inventory reservation, the compensating transaction publishes an OrderCancelled event that the Inventory Service consumes to release the reserved inventory, ensuring that inventory is not indefinitely held for orders that cannot be completed. The compensating transaction pattern ensures that partial failures do not leave the system in an inconsistent state where inventory is reserved but orders are not confirmed.

\section{Database Design}

Database design specifies the schema structures for each service's dedicated PostgreSQL database, including table definitions, column specifications, constraints, indexes, and relationships.

\subsection{User Service Database Schema}

The User Service database contains a single primary table named users storing user account information and credentials. The table includes columns for id defined as UUID type with primary key constraint serving as the unique user identifier, email defined as varchar with two hundred fifty-five character length with unique and not null constraints ensuring email uniqueness and presence, password defined as varchar with two hundred fifty-five character length with not null constraint storing the bcrypt hashed password, first\_name and last\_name defined as varchar with one hundred character length with not null constraints storing user's name components, phone\_number defined as varchar with twenty character length allowing null values for optional phone contact, role defined as varchar with fifty character length with not null constraint defaulting to USER specifying the user's authorization role, is\_active defined as boolean with not null constraint defaulting to true indicating whether the account is enabled, created\_at defined as timestamp with not null constraint defaulting to current timestamp recording account creation time, and updated\_at defined as timestamp with not null constraint defaulting to current timestamp and updating automatically on modifications recording the last update time.

Indexes are created on the email column to optimize authentication queries that look up users by email address, and on the role column to support administrative queries filtering users by role. A unique index on the email column enforces the uniqueness constraint at the database level, preventing duplicate email registrations even under concurrent registration attempts.

\subsection{Product Service Database Schema}

The Product Service database contains a products table with columns for id defined as UUID with primary key constraint, name defined as varchar with two hundred character length with not null constraint, description defined as text type allowing lengthy product descriptions, price defined as decimal with ten total digits and two decimal places with not null constraint and check constraint ensuring positive values, category defined as varchar with one hundred character length with not null constraint, stock\_quantity defined as integer with not null constraint defaulting to zero, image\_url defined as varchar with five hundred character length allowing null for products without images, is\_active defined as boolean with not null constraint defaulting to true, created\_at and updated\_at defined as timestamps with automatic management.

Full-text search support is implemented through a search\_vector column of type tsvector that combines the name, description, and category fields converted to tsvector format using the to\_tsvector function with English language configuration. A GIN index on the search\_vector column enables efficient full-text search queries, with the index automatically updated through a database trigger that recalculates the search\_vector whenever the name, description, or category columns are modified. This design provides performant search capabilities without requiring external search engines for moderate-scale product catalogs.

\subsection{Order Service Database Schema}

The Order Service database implements a normalized schema with orders table and order\_items table supporting one-to-many relationship between orders and their line items. The orders table includes columns for id defined as UUID with primary key, customer\_id defined as UUID with not null constraint referencing the User Service's user identifier, total\_amount defined as decimal with ten digits and two decimal places with not null constraint, status defined as varchar with fifty character length with not null constraint defaulting to PENDING, saga\_status defined as varchar with fifty character length tracking Saga workflow progress, created\_at and updated\_at defined as timestamps with automatic management, and version defined as integer with not null constraint defaulting to zero for optimistic locking to prevent concurrent modification conflicts.

The order\_items table includes columns for id defined as UUID with primary key, order\_id defined as UUID with not null constraint and foreign key reference to the orders table with cascade delete ensuring order items are removed when their parent order is deleted, product\_id defined as UUID with not null constraint storing the product identifier from the Product Service, product\_name defined as varchar with two hundred character length with not null constraint storing a snapshot of the product name at order time to preserve order details even if product names change, quantity defined as integer with not null constraint and check constraint ensuring positive quantities, unit\_price defined as decimal with ten digits and two decimal places with not null constraint storing the price at order time, and line\_total defined as decimal with ten digits and two decimal places calculated as quantity multiplied by unit\_price.

Indexes are created on the customer\_id column to support customer order history queries, on the status column to support order filtering by status, and on the created\_at column to support date range queries. A composite index on customer\_id and created\_at optimizes queries retrieving a customer's orders sorted by date.

\subsection{Payment Service Database Schema}

The Payment Service database contains a payments table with columns for id defined as UUID with primary key, order\_id defined as UUID with not null constraint referencing the associated order, payment\_method defined as varchar with fifty character length with not null constraint specifying the payment type such as CREDIT\_CARD, PAYPAL, or BANK\_TRANSFER, amount defined as decimal with ten digits and two decimal places with not null constraint, currency defined as varchar with three character length with not null constraint defaulting to USD, status defined as varchar with fifty character length with not null constraint defaulting to INITIATED, gateway\_transaction\_id defined as varchar with two hundred fifty-five character length allowing null to store the transaction identifier returned by the payment gateway, gateway\_response defined as JSON type storing the complete gateway response for debugging and reconciliation, created\_at and updated\_at defined as timestamps with automatic management.

\subsection{Inventory Service Database Schema}

The Inventory Service database implements an inventory table with columns for id defined as UUID with primary key, product\_id defined as UUID with not null constraint and unique constraint ensuring one inventory record per product, total\_stock defined as integer with not null constraint defaulting to zero representing total physical inventory, reserved\_stock defined as integer with not null constraint defaulting to zero representing inventory reserved for pending orders, available\_stock defined as integer with not null constraint defaulting to zero calculated as total\_stock minus reserved\_stock, low\_stock\_threshold defined as integer with not null constraint defaulting to ten specifying the threshold below which low-stock alerts are triggered, and version defined as integer with not null constraint defaulting to zero for optimistic locking.

A reservations table tracks individual inventory reservations with columns for id defined as UUID with primary key, product\_id defined as UUID with not null constraint referencing the inventory record, order\_id defined as UUID with not null constraint referencing the order that created the reservation, quantity defined as integer with not null constraint, created\_at defined as timestamp with not null constraint, expires\_at defined as timestamp with not null constraint specifying when the reservation expires if not completed, and status defined as varchar with fifty character length with not null constraint defaulting to ACTIVE. Indexes on order\_id support reservation lookup by order, and indexes on expires\_at support expiration processing queries finding expired reservations.

\section{API Design}

API design specifies the RESTful endpoints exposed by each service through the API Gateway, including request and response formats, HTTP methods, status codes, and error handling.

\subsection{User Service API Endpoints}

The User Service exposes endpoints for user registration at POST slash api slash users slash register accepting registration request body containing email, password, first name, and last name fields, validating inputs, creating the user account with USER role, and returning two hundred and one status with user details excluding password. The login endpoint at POST slash api slash users slash login accepts login request body containing email and password, validates credentials against stored hash, generates JWT token, and returns two hundred status with authentication response containing token, token type as Bearer, user ID, email, and assigned roles.

The profile retrieval endpoint at GET slash api slash users slash profile requires authentication via Bearer token, extracts user ID from token, retrieves user details from database, and returns two hundred status with user profile information. The profile update endpoint at PUT slash api slash users slash profile requires authentication, accepts updated profile fields, validates inputs, updates the database record, and returns two hundred status with updated profile. The password change endpoint at PUT slash api slash users slash password requires authentication, accepts current password and new password, verifies current password, validates new password complexity, updates the password with new bcrypt hash, and returns two hundred status with success message.

\subsection{Product Service API Endpoints}

The Product Service exposes endpoints for product retrieval at GET slash api slash products returning paginated product list with default page size of twenty, supporting query parameters for page number, page size, and sort field. The product detail endpoint at GET slash api slash products slash open curly brace id close curly brace retrieves individual product by UUID, returning four hundred and four if not found or two hundred with product details.

The product search endpoint at GET slash api slash products slash search accepts query parameter for search term, converts to tsquery, executes full-text search with ranking, and returns results sorted by relevance. The category filter endpoint at GET slash api slash products slash category slash open curly brace category close curly brace filters products by category, supporting additional query parameters for price range filtering through min\_price and max\_price parameters.

Product creation at POST slash api slash users slash products requires ADMIN role, accepts product creation request body with name, description, price, category, stock quantity, and optional image URL, validates inputs, persists to database, and returns two hundred and one with created product details. Product update at PUT slash api slash products slash open curly brace id close curly brace requires ADMIN role, accepts updated fields, validates and updates the record, and returns two hundred with updated product. Product deletion at DELETE slash api slash products slash open curly brace id close curly brace requires ADMIN role, performs soft deletion by setting is\_active to false, and returns two hundred and four no content status.

\subsection{Order Service API Endpoints}

The Order Service exposes order creation at POST slash api slash orders requiring authentication, accepting order creation request with customer ID and array of order items containing product IDs and quantities, validates product existence and availability, calculates total amount, creates order in PENDING status, publishes OrderCreated event to initiate Saga workflow, and returns two hundred and one with order details including order ID, status, total amount, and estimated processing time.

Order retrieval at GET slash api slash orders slash open curly brace id close curly brace requires authentication, validates that the requesting user owns the order or has ADMIN role, retrieves order with line items from database, and returns two hundred with complete order details. Order history at GET slash api slash orders slash customer slash open curly brace customerId close curly brace requires authentication, validates customer ID matches authenticated user or has ADMIN role, retrieves orders filtered by customer with optional status and date range filters, and returns paginated order list.

Order cancellation at POST slash api slash orders slash open curly brace id close curly brace slash cancel requires authentication, validates order ownership, checks that order is in PENDING or CONFIRMED status and not yet shipped, updates status to CANCELLED, publishes OrderCancelled event to trigger compensating transactions, and returns two hundred with cancellation confirmation.

\section{Event Schema Design}

Event schema design specifies the structure and content of Kafka events exchanged between services, ensuring consistent event formats and comprehensive information for event handlers to process events correctly.

\subsection{OrderCreated Event}

The OrderCreated event is published by the Order Service to initiate the Saga workflow, containing fields for eventId defined as UUID providing unique event identifier for deduplication, eventType defined as string with value ORDER\_CREATED, timestamp defined as ISO-8601 formatted datetime string, aggregateId defined as UUID representing the order identifier, and payload containing nested object with orderId, customerId, array of items each with productId, productName, quantity, and unitPrice, totalAmount as decimal, and createdAt as datetime string.

\subsection{StockReserved Event}

The StockReserved event is published by the Inventory Service acknowledging successful inventory reservation, containing eventId, eventType as STOCK\_RESERVED, timestamp, aggregateId as order identifier, and payload with orderId, reservationId as UUID identifying the reservation, items array with productId and reservedQuantity, reservedAt as datetime string, and expiresAt as datetime string specifying reservation expiration time.

\subsection{PaymentCompleted Event}

The PaymentCompleted event is published by the Payment Service acknowledging successful payment, containing eventId, eventType as PAYMENT\_COMPLETED, timestamp, aggregateId as order identifier, and payload with orderId, paymentId as UUID identifying the payment, amount paid, currency, paymentMethod, gatewayTransactionId, and completedAt as datetime string.

\section{Security Design}

Security design addresses authentication, authorization, data protection, and threat mitigation across the platform.

JWT-based authentication implements stateless authentication where the API Gateway validates tokens on incoming requests, extracts user information from token payload, and forwards user context to backend services through HTTP headers. Token validation verifies the HMAC-SHA256 signature using the shared secret key, checks token expiration against current time, and extracts user ID, email, and roles from the payload. Invalid or expired tokens result in four hundred and one unauthorized responses without forwarding to backend services.

Authorization enforcement occurs at two levels: the API Gateway enforces coarse-grained authorization by requiring authentication for all endpoints except login and registration, and backend services enforce fine-grained authorization by checking user roles for protected operations. The Product Service requires ADMIN role for product creation, update, and deletion endpoints while allowing all authenticated users to browse products. The Order Service requires users to own orders they are accessing, preventing customers from viewing other customers' orders.

\section{Resilience Design}

Resilience design implements fault tolerance patterns to prevent cascade failures and ensure system stability under failure conditions.

Circuit breaker configuration uses Resilience4j with failure rate threshold of fifty percent, sliding window size of ten requests, wait duration in open state of ten seconds, and permitted number of calls in half-open state of three calls. Circuit breakers are configured for inter-service calls from API Gateway to backend services and from Order Service to Product Service for price validation.

Retry configuration implements exponential backoff with initial wait duration of one second, multiplier of two for subsequent retries, maximum wait duration of ten seconds, and maximum of three retry attempts. Retries are applied only to idempotent operations such as GET requests and operations safe to repeat, avoiding retries on POST requests that may create duplicate resources.

Rate limiting uses Redis-backed token bucket algorithm with default limits of twenty requests per second for general endpoints and ten requests per second for authentication endpoints, with burst capacity of fifty requests to accommodate legitimate traffic spikes. Rate limiters are applied at the API Gateway level, rejecting excess requests with four hundred and twenty-nine Too Many Requests status and Retry-After header indicating when the client can retry.

\section{Summary}

This chapter has presented comprehensive system design covering high-level architecture with service decomposition strategy, component design with layered architecture pattern and Saga implementation, database design with detailed schema specifications for each service, API design with RESTful endpoint specifications, event schema design for Kafka messaging, security design with JWT authentication and role-based authorization, and resilience design with circuit breakers, retries, and rate limiting.

The design satisfies all functional and non-functional requirements specified in the previous chapter, implementing microservices architecture patterns, event-driven communication, distributed transaction management, and comprehensive security controls. The design serves as the blueprint for the implementation detailed in the following chapters, with clear mapping from design decisions to implementation artifacts.
